\appendixsection{Листинги исходного кода}

В данном приложении приведены фрагменты исходного кода основных компонентов системы.

\subsubsection{Создание приложения и маршрутизация (Rust)}

Ниже приведён фрагмент файла \texttt{backend/src/lib.rs}, отвечающий за создание экземпляра приложения Axum с маршрутизацией и middleware.

\begin{lstlisting}[language=C,basicstyle=\small\ttfamily]
pub fn create_app(state: AppState) -> Router {
    let api_routes = Router::new()
        .nest("/auth",       handlers::auth::router())
        .nest("/animals",    handlers::animals::router())
        .nest("/milk",       handlers::milk::router())
        .nest("/reproduction",
                                handlers::reproduction::router())
        .nest("/feed",       handlers::feed::router())
        .nest("/fitness",    handlers::fitness::router())
        .nest("/reports",    handlers::reports::router())
        .nest("/analytics",  handlers::analytics::router())
        .nest("/settings",   handlers::settings::router())
        .nest("/lely",       handlers::lely::router());

    Router::new()
        .nest("/api/v1", api_routes)
        .layer(CsrfLayer)
        .layer(RateLimitLayer)
        .layer(MetricsLayer)
        .layer(CompressionLayer)
        .layer(CorsLayer)
        .with_state(state)
}
\end{lstlisting}

\subsubsection{Экстрактор авторизации (Rust)}

Ниже приведена реализация экстрактора \texttt{Claims}, извлекающего и проверяющего JWT-токен из запроса.

\begin{lstlisting}[language=C,basicstyle=\small\ttfamily]
pub struct Claims {
    pub sub: String,
    pub role: String,
    pub must_change_password: bool,
    pub jti: String,
    pub exp: i64,
}

#[async_trait]
impl<S: Send + Sync> FromRequestParts<S> for Claims {
    type Rejection = AppError;

    async fn from_request_parts(
        parts: &mut Parts,
        _state: &S,
    ) -> Result<Self, Self::Rejection> {
        let token = extract_token(parts)?;
        let claims = decode_jwt(&token)?;
        if claims.must_change_password {
            return Err(AppError::Forbidden);
        }
        check_revoked(&claims.jti).await?;
        Ok(claims)
    }
}
\end{lstlisting}

\subsubsection{Обучение модели XGBoost (Python)}

Ниже приведён фрагмент сервиса обучения модели прогнозирования надоев.

\begin{lstlisting}[language=Python,basicstyle=\small\ttfamily]
def train_milk_forecast(df: pd.DataFrame) -> dict:
    features = [
        "milk_7d_avg", "milk_14d_avg", "milk_30d_avg",
        "milk_7d_std", "milk_14d_std",
        "dim", "lactation_number", "fpr",
        "month", "day_of_week",
    ]
    X = df[features]
    y = df["milk_kg"]

    models = {}
    for q in [0.1, 0.5, 0.9]:
        model = xgb.XGBRegressor(
            objective="reg:quantileerror",
            quantile_alpha=q,
            n_estimators=200,
            max_depth=6,
            learning_rate=0.05,
            subsample=0.8,
            colsample_bytree=0.8,
        )
        model.fit(X, y)
        models[q] = model

    return models
\end{lstlisting}

\subsubsection{API-клиент (TypeScript)}

Ниже приведён фрагмент типизированного API-клиента на стороне frontend.

\begin{lstlisting}[language=Java,basicstyle=\small\ttfamily]
export const animalsApi = {
  list: (params: AnimalFilter) =>
    get<PaginatedResponse<Animal>>(
      "/animals?" + buildQuery(params)
    ),
  getById: (id: number) =>
    get<Animal>(`/animals/${id}`),
  create: (data: CreateAnimal) =>
    post<Animal>("/animals", data),
  update: (id: number, data: UpdateAnimal) =>
    put<Animal>(`/animals/${id}`, data),
  deactivate: (id: number) =>
    del(`/animals/${id}`),
  batchDeactivate: (ids: number[]) =>
    post("/animals/batch/deactivate", { ids }),
};
\end{lstlisting}
